\begin{summarybox}
\begin{minipage}{\linewidth}
\textbf{一句话：先求圆心距离范围，再看半径落在范围的哪里。}
\[
 d_{\min}=\min_{P\in E}PM,\qquad d_{\max}=\max_{P\in E}PM,
 \qquad D_{\max}=d_{\max}+r.
\]
\[
 D_{\min}=
 \begin{cases}
 d_{\min}-r,&r<d_{\min},\\
 0,&d_{\min}\le r\le d_{\max},\\
 r-d_{\max},&r>d_{\max}.
 \end{cases}
\]
\textbf{按图核对：}
\begin{itemize}[leftmargin=*,itemsep=3pt]
\item 外离或小圆在椭圆内：半径小于所有 $PM$，最近距离取 $d_{\min}-r$。
\item 相交或边界相切：存在公共点，最近距离为0。
\item 圆包住椭圆：半径大于所有 $PM$，最近距离取 $r-d_{\max}$。
\item 最远距离：椭圆点、圆心、圆周点依次共线，始终取和。
\end{itemize}
\textbf{最后检查：}研究的是圆周还是圆盘？是否求出了完整距离范围？
是否计入相切边界？图上的极值点能否代回两个方程？
\end{minipage}
\end{summarybox}
